Nel corso della Storia dell'Informatica, l'avanzamento tecnico e tecnologico verificatosi nell'ambito delle Architetture dei Calcolatori e dei Sistemi Operativi ha permesso grandi passi in avanti nelle dinamiche di utilizzo degli Elaboratori Elettronici. Si è assistito al passaggio dall'esecuzione sequenziale delle operazioni ad un'esecuzione parallela delle stesse, il tutto all'interno di un singolo Elaboratore. \par
L'opportunità di avere esperienze utente multiprogrammate è stata resa possibile, oltre che dagli Elaboratori non sequenziali, anche da Linguaggi di Programmazione anch'essi non sequenziali che permettano l'esecuzione concorrente di più processi. Tale flessibilità ha permesso un grado di utilizzo delle risorse nettamente superiore al passato, garantendo prestazioni superiori a parità di risorse. \par
È chiaro come il Linguaggio di Programmazione Java sia stato nel tempo un'ottima piattaforma per lo sviluppo di Software che sfrutti appieno queste possibilità, contribuendo significativamente allo sviluppo dell'Informatica come la conosciamo oggi.